%%
%% Ergänzung zu Kapitel 5: native iOS-/iPadOS-Implementierung
%%

\begin{neu}
\section{Native iOS-/iPadOS-Implementierung}

Die Apple-Portierung wurde als eigenständige native Anwendung umgesetzt. Die oberste Darstellungsschicht besteht aus SwiftUI-Views, deren Zustände durch ein zentrales App-Modell bereitgestellt werden. Fachliche Regeln liegen in einem logisch abgegrenzten Pure-Swift-Domainbereich. Kritische Abläufe wie Aktivierung, Persistenz und Netzwerkkoordination werden über serialisierte Koordinatoren gekapselt. Dadurch greifen Views nicht direkt auf Keychain, Dateien oder Netzwerkservices zu. Die wesentliche Abhängigkeitsrichtung lässt sich vereinfacht als

\[
\text{SwiftUI} \rightarrow \text{AppModel} \rightarrow \text{Domain/Coordinator} \rightarrow \text{Infrastructure}
\]

beschreiben. Diese Struktur bildet dieselbe fachliche Trennung wie die Android-Anwendung ab, nutzt jedoch die nativen Apple-Frameworks.

Der Funktionsumfang umfasst den vollständigen für die Studie vorgesehenen App-Ablauf: deutsch- und englischsprachige Oberfläche, Informationsfeed, lokale Ausblendung gelesener Informationen, Benachrichtigungsentscheidung, optionale lokale Erinnerungen, Aktivierung über E-Mail-Adresse oder Prolific-ID, lokale Demo, Feedback, serverseitige Feature-Freigabe, produktive Matching- und Labeling-Durchläufe, Craving-Selbstbericht, Wiederholungslogik für ausstehende Übertragungen sowie die beiden Abschlusswege \texttt{COMPENSATION\_CODE} und \texttt{PROLIFIC\_MANUAL}. Die Implementierung verwendet dabei die bestehenden Backendverträge und führt für iOS keine eigene wissenschaftliche API ein.

Die Aktivierung folgt demselben zweistufigen Protokoll wie unter Android. Eine syntaktisch gültige E-Mail-Adresse oder eine 24-stellige alphanumerische Prolific-ID wird an den bestehenden Aktivierungsendpunkt übergeben. Das vom ersten Request erhaltene UUID-v4-App-Token bleibt zunächst flüchtig. Erst wenn der zweite Request mit Identifier und Token erfolgreich mit HTTP~204 bestätigt wurde, wird das Token dauerhaft gespeichert. Ein mehrdeutiger Abbruch während dieser Bestätigung wird nicht automatisch als erfolgreiche Aktivierung interpretiert. Damit bleibt die serverseitig definierte Aktivierungssemantik plattformübergreifend erhalten.

Der produktive Studienablauf wird aus den kanonischen Manifesten erzeugt. Für Matching werden alle 50 vorgesehenen Items genau einmal einer zufälligen Reihenfolge zugeordnet und anschließend in zehn Durchläufe zu je fünf Trials zerlegt. Für Labeling werden die festgelegten Blöcke verwendet. Die konkrete Auswahlposition wird pro Trial zufällig angeordnet. Nach fünf Trials folgt die Rauchverlangensabfrage. Ein Prozessabbruch vor dem Absenden eines Selbstberichts persistiert weder Trialposition noch konkrete Auswahl; der noch nicht bestätigte Durchlauf beginnt nach Wiederaufnahme erneut. Damit werden nur Zustände dauerhaft gespeichert, die für Studienfortschritt und technische Recovery notwendig sind.

Vor einem Submission-Request persistiert die App ausschließlich den ganzzahligen Craving-Wert als ausstehend. Der anschließende Request enthält das App-Token, den Craving-Wert und die unveränderte App-Version. Situation, Versuchsbedingung, konkrete Auswahl, Sprache und Plattform werden nicht vom Client übertragen. Die Zuordnung zu Matching oder Labeling bleibt somit wie bei Android serverseitig vom bereits gespeicherten Fortschritt abhängig. Erst nach einer syntaktisch und semantisch gültigen Serverantwort wird der lokale Fortschritt atomar aktualisiert. Ein ausstehender Wert sperrt bis zur erfolgreichen Wiederholung den Beginn einer weiteren Studiensituation.

Auch die Abschlusslogik entspricht den bereits beschriebenen beiden Backendmodi. Beim direkten Abschluss wird der empfangene UUID-v4-Abrechnungscode zunächst als noch zu bestätigender Zustand gespeichert. Anschließend erfolgt der separate Bestätigungsrequest; bei dessen Fehlschlag bleibt der Code für einen späteren Retry erhalten und wird der teilnehmenden Person noch nicht als vollständig bestätigter Abschluss präsentiert. Der Prolific-Pfad speichert keinen CueLens-Abrechnungscode, sondern übernimmt den vom Backend gelieferten Abschlussmodus \texttt{PROLIFIC\_MANUAL}. Damit bleibt die administrative Unterscheidung der Rekrutierungskanäle erhalten, ohne die wissenschaftliche Datenstruktur plattformspezifisch zu verändern.

Die Apple-Anwendung trennt außerdem unkritische UI-Einstellungen von studienrelevanter Persistenz. Spracheinstellungen und positive IDs bereits bekannter beziehungsweise dauerhaft ausgeblendeter Feednachrichten liegen in \texttt{UserDefaults}. App-Token, Aktivierungs-Recovery und Studienzustand werden dagegen in hierfür vorgesehenen geschützten Speichern abgelegt. Nachrichteninhalte, Trialentscheidungen und Reaktionszeiten werden nicht dauerhaft gespeichert. Diese Trennung reduziert die Menge persistent gehaltener Daten und verhindert, dass fachlich sensitive Zustände allein über ungeschützte Komfortspeicher verwaltet werden.

Für Hintergrundfunktionen wurde bewusst keine Abhängigkeit der Kernstudie von iOS-Scheduling angenommen. Lokale Benachrichtigungen enthalten keine sensiblen Nutzdaten, keinen Ton und keinen Badge. Studienerinnerungen werden aus dem aktuellen fachlichen Zustand abgeleitet und bei nicht mehr zulässigem Zustand entfernt. Die Hintergrundaktualisierung des Informationsfeeds ist lediglich best effort; schlägt sie aus oder wird sie vom Betriebssystem nicht eingeplant, bleiben Aktivierung, Studiendurchführung und Submission weiterhin funktionsfähig.

Die Portierung ist damit als eigenständige vollständige Clientimplementierung zu bewerten, nicht als bloßer UI-Prototyp. Gleichzeitig ist zwischen Implementierungsreife und Distributionsreife zu unterscheiden. Das interne Definition-of-Done-Audit weist die Architektur, den funktionalen Umfang, die Studienmethodik und die automatisierbaren Sicherheits- und Testkriterien als implementiert beziehungsweise bestanden aus. Die formale Freigabe bleibt jedoch blockiert, solange reale Staging-Ende-zu-Ende-Prüfungen, Distribution-Signing, TestFlight-Prüfung und externe Freigaben nicht abgeschlossen sind. Für die Masterarbeit ist diese Differenz relevant, weil sie einen technisch vollständigen Prototypen belegt, ohne daraus eine bereits produktiv veröffentlichte iOS-App abzuleiten.
\end{neu}
